% CV publico del portafolio (ES). Copia de empleo-agente/salidas/aplicaciones/_cv_base/es/cv.tex.
% Sin nombre de la empresa, sin telefono y sin el abrepuertas (es publico en la web).
% Sin guiones largos (---) a proposito; "--" solo en rangos de fechas.
\documentclass[a4paper,10pt]{article}
\usepackage[top=0.6in, bottom=0.6in, left=0.7in, right=0.7in]{geometry}
\usepackage[T1]{fontenc}
\usepackage[utf8]{inputenc}
\usepackage[spanish]{babel}
\usepackage{lmodern}
\usepackage{microtype}
\usepackage{enumitem}
\usepackage{needspace}
\usepackage{parskip}
\usepackage{titlesec}
\usepackage[dvipsnames]{xcolor}
\definecolor{acc}{HTML}{1F3A5F}
\usepackage[colorlinks=true,urlcolor=acc,linkcolor=acc,bookmarks=false,
  pdfauthor={Juan Nicolás García Vega},pdftitle={CV Juan Nicolás García Vega}]{hyperref}

\pagestyle{empty}
\raggedright
\hyphenpenalty=10000 \exhyphenpenalty=10000
\emergencystretch=3em
\clubpenalty=10000 \widowpenalty=10000

\setlist[itemize,1]{leftmargin=1em,   itemsep=3pt,   parsep=0pt, topsep=2pt,   label=\textbullet}
\setlist[itemize,2]{leftmargin=1.4em, itemsep=1.5pt, parsep=0pt, topsep=1.5pt, label=\textendash}
\titleformat{\section}{\large\bfseries\uppercase}{}{0em}{}[\vspace{0.1em}\titlerule]
\titlespacing{\section}{0pt}{0.8em}{0.5em}
\newcommand{\entry}[3]{\textbf{#1}\ifx&#2&\else, \textit{#2}\fi \hfill \mbox{\textit{#3}}}
\newcommand{\gh}[1]{\href{https://github.com/Nicolas6879/#1}{github.com/Nicolas6879/#1}}

\begin{document}

\begin{center}
    {\huge \textbf{Juan Nicolás García Vega}} \\[0.3em]
    Ingeniero Full-Stack \textbar{} Vue, React, Laravel, Node.js y Python \\[0.35em]
    Bucaramanga, Colombia \textperiodcentered{} Disponible para reubicación o remoto \\
    \href{mailto:jngv08@gmail.com}{jngv08@gmail.com} \textperiodcentered{}
    \href{https://www.linkedin.com/in/juan-nicolas-garcia-vega}{linkedin.com/in/juan-nicolas-garcia-vega} \textperiodcentered{}
    \href{https://github.com/Nicolas6879}{github.com/Nicolas6879}
\end{center}

\needspace{6\baselineskip}
\section*{Perfil}
Ingeniero full-stack (Ingeniería de Sistemas, julio de 2026) que ha llevado
\textbf{8 sistemas a producción} en una empresa industrial colombiana con Vue 3, TypeScript, Laravel y SQL Server, entre
ellos un \textbf{sistema de gestión de calidad ISO 9001 que reemplazó más de 600 documentos en
papel}. Construye aplicaciones web de extremo a extremo, APIs REST y herramientas con IA (MCP,
un plugin para Claude Code), con pruebas automáticas. \textbf{Hedera Developer \& Community
Ambassador en Colombia}; \textbf{primer lugar} en el CTF Web3 de Hedera DevDay y \textbf{segundo
lugar} en el bounty de automatización on-chain, ETHDenver 2026.

\needspace{6\baselineskip}
\section*{Experiencia profesional}
\begin{itemize}
    \item \entry{Ingeniero de Software}{Empresa industrial, Colombia}{Abr.\ 2026 -- Presente}
        \begin{itemize}
            \item Lidero de punta a punta la transición a \textbf{operaciones sin papel}: diseño
                y construyo los sistemas, entre ellos la aplicación de expedientes digitales
                (PDFs escaneados, archivados y firmados) y, además, reestructuré y completé el
                sistema anterior de adjuntos (PDFs generados por el ERP, antes manejados a mano).
            \item Construí un \textbf{servidor MCP} para buscar en los catálogos del SGC
                (Sistema de Gestión de Calidad) y \textbf{publicar documentos de calidad desde
                Claude}, para que la gente deje de saltarse el paso de subirlos.
            \item Mantengo y extiendo el \textbf{ERP} (Visual Basic, SQL Server), que usan unas
                \textbf{100 personas}.
            \item Hago parte del equipo que construye un \textbf{asistente de IA} interno
                (Gemini API, migrando a Claude) para el conocimiento de la empresa y los reportes.
            \item Entregué \textbf{Preoperacionales} (listas de chequeo de equipos en 3 sedes,
                modo sin conexión con sincronización automática), \textbf{ComidApp} (pedidos y
                costeo multiplanta), un portal de tableros \textbf{Power BI}, un formulario de
                debida diligencia de terceros para \textbf{SAGRILAFT} (normas colombianas contra
                el lavado de activos) con salida en PDF, y la migración completa de un sistema
                en \textbf{C\#} a Vue 3, Laravel y SQL Server.
        \end{itemize}

    \item \entry{Practicante de Ingeniería de Sistemas}{Empresa industrial, Colombia}{Ago.\ 2025 -- Abr.\ 2026}
        \begin{itemize}
            \item Construí desde cero el \textbf{Sistema de Gestión de Calidad ISO 9001}
                (trabajo de grado), de uso diario en \textbf{3 áreas}, que reemplazó
                \textbf{más de 600 documentos en papel}. Seis módulos: documentos, indicadores,
                auditorías, riesgos, acciones CAPA y no conformidades, gestión del cambio.
                Vue 3, TypeScript, Pinia, Laravel, MySQL.
        \end{itemize}

    \item \entry{Developer \& Community Ambassador, Colombia}{Hedera}{2025 -- Presente}
        \begin{itemize}
            \item Organicé, dicté y fui jurado de \textit{Descifra Hedera: Charla y Hackathon
                Web3} (29 sep.\ 2026, U26 Fest, Universidad Industrial de Santander): 3 horas,
                unas \textbf{90 personas}.
            \item Dicté dos talleres virtuales: el arranque local del \textbf{Hedera APEX
                Hackathon} para México (26 feb.\ 2026, 3 h) y la incorporación de estudiantes de
                dos colegios de comunidades vulnerables de Brasil al Hedera Developer Portal, para
                construir su primera dapp (27 feb.\ 2026, 3,5 h, dentro de un programa de inclusión).
            \item Speaker en \textbf{Blockchain Summit Latam Bogot\'a} (4 nov.\ 2026): taller de
                1 hora \textit{Tokenización de activos y agentes de IA con control y trazabilidad
                en Hedera}, con el embajador Juan Manuel G\'omez.
            \item Fui jurado por encargo de Hedera en ETHOnline 2026 (unos \textbf{200 proyectos})
                y NEA Tech 2026, Argentina (unos \textbf{20 proyectos}).
            \item Acompaño a desarrolladores nuevos y apoyo a equipos con contratos inteligentes y
                SDKs; representé a Hedera en ETHDenver 2026, Hedera DevDay y Consensus Miami 2026.
        \end{itemize}

    \item \entry{Líder técnico}{X-Culture Professional Association (XPA)}{2025 -- 2026}
        \begin{itemize}
            \item Construí el sitio oficial (WordPress, Elementor) con pasarela de pagos y llevé
                la parte técnica de eventos virtuales internacionales.
        \end{itemize}

    \item \entry{Instructor y mentor de programación}{Kodland}{2021 -- 2024}
        \begin{itemize}
            \item Enseñé \textbf{Python} a más de \textbf{150 estudiantes} (de 10 a 20 años) en más de 30 cursos.
        \end{itemize}
\end{itemize}

\needspace{6\baselineskip}
\section*{Proyectos destacados}
\begin{itemize}
    \needspace{4\baselineskip}
    \item \textbf{Tacaño} (\gh{tacano}), plugin de código abierto para Claude Code: reduce el
        uso de tokens en cerca de \textbf{39\,\%}, medido reproduciendo \textbf{500 turnos
        reales}. Un hook consulta a un modelo pequeño cuándo compactar y qué delegar a
        trabajadores más baratos. En un catálogo de plugins, con documentación en inglés y español.
    \needspace{4\baselineskip}
    \item \textbf{Willys por Colombia} (\href{https://willys-por-colombia.vercel.app}{demo},
        \gh{willys-por-colombia}): juego de manejo en three.js con música en vivo de
        \textbf{Google Lyria RealTime}. La música se mezcla según la cercanía a cada región y se
        va desviando al manejar; cada región repite su tema desde una semilla fija. 18 suites de
        pruebas automáticas y controles táctiles.
    \needspace{4\baselineskip}
    \item \textbf{VerifiChain Visual} (\gh{verifichain-vision-oracle}), hackathon Hedera APEX
        2026: agente de inspección visual para metal y concreto (YOLOv8 con respaldo de Claude).
        Cada veredicto queda registrado en Hedera y certificado con un NFT. Pensado para
        empresas industriales de Santander.
    \needspace{4\baselineskip}
    \item \textbf{x402 Guard} (\gh{x402-guard}), Hedera x402 Bounty: proxy de pagos para agentes
        de IA. El agente no tiene llaves; el proxy paga solo si una política de gasto lo permite,
        y cada decisión, aprobada o bloqueada, queda en un registro de auditoría público (Hedera
        Consensus Service).
    \needspace{4\baselineskip}
    \item \textbf{Hive Protocol} (\gh{HIVE}), ETHDenver 2026, \textbf{segundo lugar} en el bounty
        de automatización on-chain de Hedera: tres agentes de IA proponen, ejecutan y validan
        conocimiento sin intervención humana, con un orquestador MCP y temas on-chain.
    \item \textbf{Más:} proyectos de hackathon ScoutBrief, Quorum Treasury, PULSO Tanda y Roast
        Pitch (\href{https://github.com/Nicolas6879}{github.com/Nicolas6879}).
\end{itemize}

\needspace{6\baselineskip}
\section*{Educación}
\begin{itemize}
    \item \entry{Universidad Industrial de Santander (UIS)}{Ingeniería de Sistemas}{2021 -- jul.\ 2026}
    \item \entry{SENA}{Técnico en Equipos Electrónicos Industriales}{2018 -- 2019}
\end{itemize}

\needspace{6\baselineskip}
\section*{Habilidades técnicas}
\begin{itemize}
    \item \textbf{Lenguajes:} Python, TypeScript, JavaScript, PHP, Visual Basic, SQL, Solidity
        (también Rust, Java, C++, Go)
    \item \textbf{Web:} Vue 3, React, Next.js, three.js, Laravel, Node.js, REST
    \item \textbf{IA y agentes:} MCP, plugins y hooks de Claude Code, llamada a herramientas,
        sistemas multiagente, APIs de LLMs (Claude, Gemini, OpenAI), YOLOv8
    \item \textbf{Blockchain:} Hedera (servicio de consenso, tokens, contratos inteligentes), Foundry, Hardhat
    \item \textbf{Datos y herramientas:} SQL Server, MySQL, PostgreSQL, Power BI, Docker, Git, Vercel, n8n
    \item \textbf{Ingeniería inversa:} análisis de tráfico, binarios de Android y Flutter, REA
\end{itemize}

\section*{Idiomas}
Español (nativo) \textperiodcentered{} Inglés (intermedio, B1; mejorando hacia B2)

\end{document}
